% Bagian: first-order-logic
% Bab: beyond
% Subbagian: introduction

\documentclass[../../../include/open-logic-section]{subfiles}

\begin{document}

\olfileid[id]{fol}{byd}{int}

\olsection{Ikhtisar}

Logika orde pertama bukan satu-satunya sistem logika yang menarik: terdapat
banyak perluasan dan variasi logika orde pertama. Suatu logika biasanya
terdiri atas spesifikasi formal sebuah bahasa, biasanya---meskipun tidak
selalu---suatu sistem deduktif, dan biasanya---meskipun tidak selalu---suatu
semantik yang dimaksudkan. Namun, penggunaan teknis istilah tersebut
memunculkan pertanyaan yang jelas: apa hubungan antara berbagai logika yang
bukan logika orde pertama dengan kata ``logika'' dalam arti intuitif atau
filosofis? Semua sistem yang diuraikan di bawah ini dirancang untuk memodelkan
suatu bentuk penalaran; dapatkah kita mengatakan apa yang menjadikannya logis?

Tidak ada jawaban mudah yang segera tersedia. Kata ``logika'' digunakan
dengan cara yang berbeda-beda dan dalam konteks yang berbeda-beda, dan
gagasannya, seperti gagasan ``kebenaran'', telah dianalisis dari berbagai
sudut pandang filosofis. Misalnya, orang dapat menganggap tujuan penalaran
logis sebagai menentukan pernyataan mana yang niscaya benar, benar secara
a priori, benar tanpa bergantung pada interpretasi istilah-istilah
nonlogis, benar berdasarkan bentuknya, atau benar berdasarkan konvensi
linguistik; dan setiap konsepsi ini memerlukan banyak penjernihan. Bahkan jika
perhatian dibatasi pada jenis logika yang digunakan dalam matematika, hanya
ada sedikit kesepakatan mengenai cakupannya. Misalnya, dalam
\textit{Principia Mathematica}, Russell dan Whitehead berusaha mengembangkan
matematika di atas landasan logika, mengikuti tradisi {\em logisisme} yang
dirintis oleh Frege. Sistem logika mereka merupakan suatu bentuk logika tipe
tinggi yang mirip dengan sistem yang diuraikan di bawah ini. Pada akhirnya
mereka terpaksa memperkenalkan aksioma-aksioma yang, menurut kebanyakan
patokan, tidak tampak murni logis (terutama aksioma ketakhinggaan dan aksioma
reduksibilitas), tetapi orang tetap dapat berpendapat bahwa beberapa bentuk
penalaran orde tinggi patut diterima sebagai penalaran logis. Sebaliknya,
Quine, yang ontologinya tidak menerima ``proposisi'' sebagai objek wacana yang
sah, berpendapat bahwa logika orde kedua dan orde tinggi sebenarnya merupakan
teori himpunan yang menyamar; dengan kata lain, sistem yang melibatkan
kuantifikasi atas predikat tidak murni logis.

Untuk saat ini, sebaiknya persoalan-persoalan filosofis semacam itu disisihkan
untuk lain waktu, dan sistem-sistem di bawah ini dipandang saja sebagai
idealisasi formal berbagai jenis penalaran, baik logis maupun tidak.

\end{document}
